\scanpage{190}
% archive sheet number: 189

Soit $\underline C$ une catégorie, \unsure{déf:|d'où}
\[
\operatorname{Pro}(\underline C)\unsure{\subset|=}\underline{\mathrm{Hom}}(\underline C,(\mathrm{Ens}))
\]
Un objet $\unsure{\mathscr X|\mathcal X}$ de $\operatorname{Pro}(\underline C)$ est connu (à isomorphisme près) quand on connait le foncteur $\mathscr Y\mapsto\mathrm{Hom}_{\operatorname{Pro}(\underline C)}(\mathscr Y,\mathscr X)$ de $\operatorname{Pro}(\underline C)^\circ$ dans $(\mathrm{Ens})$, \deleted{\unsure{mais}} mais,

\deleted{$\mathscr Y=(Y_j)_{j\in\unsure{I|J}}$, les $\unsure{Y_i|Y_j}\in\underline C$, alors \unsure{...},}
\[
\deleted{\mathrm{Hom}_{\operatorname{Pro}(\underline C)}(\mathscr Y,\mathscr X)=\varprojlim_{\unsure{i|j}}\mathrm{Hom}(\mathscr Y,\mathscr X_i)=}
\]
identifiant $\underline C$ à une sous-catégorie de $\operatorname{Pro}(\underline C)$, on voit que $\mathscr X$ n'est \uline{plus connu} à isom près quand on connait le foncteur $Y\mapsto\mathrm{Hom}_{\operatorname{Pro}(\underline C)}(Y,\mathscr X)$ de $\underline C^\circ$ dans $(\mathrm{Ens})$. [Ex: $\underline C=(\mathrm{Ens})$, alors $(X_i)$ et $X=\varprojlim_{i\,(\mathrm{Ens})}X_i$ définissent le $=$ foncteur $\underline C^\circ\mathbin{-}(\mathrm{Ens})$, mais en général ne sont pas isomorphes ds $\operatorname{Pro}(\underline C)$]. Le foncteur naturel
\[
\operatorname{Pro}(\underline C)\longrightarrow\underline{\mathrm{Hom}}(\underline C^\circ,(\mathrm{Ens}))
\]
n'est \inserted{donc} \uline{pas} pleinement fidèle en général. Il importe de bien distinguer en particulier, si $X=\varprojlim_{\underline C}\mathscr X$ existe, entre $\mathscr X$ et $X$

[on peut dire seulement qu'on a un morphisme naturel ds $\operatorname{Pro}\underline C$
\[
X=\varprojlim_{\underline C}\mathscr X\longrightarrow\mathscr X
\]
